% GENERATED FILE. Edit canonical lessons, then run npm run generate:books.
% canonical-source-hash: fnv1a64:6ca844d7386c5311
% canonical-lessons: PA-C178-vaje, PA-C178-ton, PA-C178-tak, PA-C178-naal, PA-C178-baad

\chapter{From Eight Until Ten, Then Home}
\label{ch:pa-a2-schedule-sequence}
\hlchaptermodality{178}

\begin{chapteropening}
\textbf{By the end of this chapter:} \emph{I can read and give one simple schedule with a whole-hour time, a source, an end limit, a companion, and one after-time step.}

Retrieve one two-sentence schedule note through listening, speaking, reading, and bounded Gurmukhi writing after building five high-reuse forms one joint at a time.
\end{chapteropening}

\section[\texorpdfstring{\pa{ਅੱਠ} \pa{ਵਜੇ} (atth vaje)}{ਅੱਠ ਵਜੇ (atth vaje)}]{\pa{ਅੱਠ} \pa{ਵਜੇ} — one familiar number becomes a clock time}

\label{lesson:PA-C178-vaje}



\begin{quote}
Say the familiar phrase \textbf{\pa{ਇਹ} \pa{ਬਹੁਤ} \pa{ਚੰਗੀ} \pa{ਕੁੜੀ} \pa{ਹੈ}} once. Then say the familiar number \textbf{\pa{ਅੱਠ}}. This lesson keeps that number and adds one time word after it.
\end{quote}

\begin{grammarlens}[title={Grammar Lens: \pa{ਵਜੇ} follows the whole hour}]
\begin{quote}
\textbf{\pa{ਅੱਠ} \pa{ਵਜੇ}} — \emph{aṭṭh vaje} — at eight o'clock
\end{quote}

Keep the order number + \textbf{\pa{ਵਜੇ}}. Here \textbf{\pa{ਵਜੇ}} marks the whole-hour time. This first step does not add half hours, minutes, or a general clock paradigm.
\end{grammarlens}

\subsection*{Your turn}
\begin{itemize}
\raggedright
  \item \emph{Hear:} \emph{aṭṭh vaje}; YOU SAY: \textbf{\pa{ਅੱਠ} \pa{ਵਜੇ}}
  \item \emph{Say it:} \textbf{\pa{ਅੱਠ} \pa{ਵਜੇ}}
  \item \emph{Say it:} the whole phrase once more — \textbf{\pa{ਅੱਠ} \pa{ਵਜੇ}}
\end{itemize}

\begin{tcolorbox}[breakable,colback=teal!4,colframe=teal!35!black,title={Before you move on}]
Say \emph{at eight o'clock}. (\textbf{\pa{ਅੱਠ} \pa{ਵਜੇ}}.) Which piece follows the number? (\textbf{\pa{ਵਜੇ}}.)

Source: Ranbir Johal, \emph{Speaking and Writing Punjabi} (2018), pp. 69–70, \href{https://openlibrary-repo.ecampusontario.ca/jspui/bitstream/123456789/492/2/2018Pun1100CoursePackage.pdf}{whole-hour vocabulary and examples}.
\end{tcolorbox}
\section[\texorpdfstring{\pa{ਘਰ} \pa{ਤੋਂ} \pa{ਸਕੂਲ} \pa{ਜਾਣਾ} (ghar ton sakul jana)}{ਘਰ ਤੋਂ ਸਕੂਲ ਜਾਣਾ (ghar ton sakul jana)}]{\pa{ਘਰ} \pa{ਤੋਂ} \pa{ਸਕੂਲ} \pa{ਜਾਣਾ} — name where a route starts}

\label{lesson:PA-C178-ton}



\begin{quote}
First retrieve \textbf{\pa{ਇਹ} \pa{ਬਹੁਤ} \pa{ਚੰਗੀ} \pa{ਕੁੜੀ} \pa{ਹੈ}} once. Then recall \textbf{\pa{ਘਰ}}, \textbf{\pa{ਸਕੂਲ}}, and \textbf{\pa{ਜਾਣਾ}}. Keep the route in that meaning order: home, then school, then go.
\end{quote}

\begin{grammarlens}[title={Grammar Lens: \pa{ਤੋਂ} follows the source}]
\begin{quote}
\textbf{\pa{ਘਰ} \pa{ਤੋਂ} \pa{ਸਕੂਲ} \pa{ਜਾਣਾ}} — \emph{ghar tõ sakūl jāṇā} — to go from home to school
\end{quote}

Punjabi uses a postposition: \textbf{\pa{ਤੋਂ}} comes after its noun. In this bounded route, \textbf{\pa{ਘਰ}} is the source, so the pair is \textbf{\pa{ਘਰ} \pa{ਤੋਂ}}. Do not move \textbf{\pa{ਤੋਂ}} before \textbf{\pa{ਘਰ}} as English moves \emph{from}.
\end{grammarlens}

\subsection*{Your turn}
\begin{itemize}
\raggedright
  \item \emph{Hear:} \textbf{\pa{ਘਰ} \pa{ਤੋਂ}}; YOU SAY: the route starts at home
  \item \emph{Say it:} \textbf{\pa{ਘਰ} \pa{ਤੋਂ} \pa{ਸਕੂਲ} \pa{ਜਾਣਾ}}
  \item \emph{Read it:} \textbf{\pa{ਘਰ} \pa{ਤੋਂ} \pa{ਸਕੂਲ} \pa{ਜਾਣਾ}}
\end{itemize}

\begin{tcolorbox}[breakable,colback=teal!4,colframe=teal!35!black,title={Before you move on}]
Build \emph{to go from home to school}. (\textbf{\pa{ਘਰ} \pa{ਤੋਂ} \pa{ਸਕੂਲ} \pa{ਜਾਣਾ}}.) Which word comes after the source noun? (\textbf{\pa{ਤੋਂ}}.)

Source: Ranbir Johal, \emph{Speaking and Writing Punjabi} (2018), pp. 67–68, \href{https://openlibrary-repo.ecampusontario.ca/jspui/bitstream/123456789/492/2/2018Pun1100CoursePackage.pdf}{from-until schedule frame}.
\end{tcolorbox}
\section[\texorpdfstring{\pa{ਅੱਠ} \pa{ਵਜੇ} \pa{ਤੋਂ} \pa{ਦਸ} \pa{ਵਜੇ} \pa{ਤੱਕ} (atth vaje ton …)}{ਅੱਠ ਵਜੇ ਤੋਂ ਦਸ ਵਜੇ ਤੱਕ (atth vaje ton …)}]{\pa{ਅੱਠ} \pa{ਵਜੇ} \pa{ਤੋਂ} \pa{ਦਸ} \pa{ਵਜੇ} \pa{ਤੱਕ} — give both ends of one time span}

\label{lesson:PA-C178-tak}



\begin{quote}
Say \textbf{\pa{ਅੱਠ} \pa{ਵਜੇ}}, then \textbf{\pa{ਦਸ} \pa{ਵਜੇ}}. Put familiar \textbf{\pa{ਤੋਂ}} after the starting time: \textbf{\pa{ਅੱਠ} \pa{ਵਜੇ} \pa{ਤੋਂ}}.
\end{quote}

\begin{grammarlens}[title={Grammar Lens: \pa{ਤੱਕ} follows the end limit}]
\begin{quote}
\textbf{\pa{ਅੱਠ} \pa{ਵਜੇ} \pa{ਤੋਂ} \pa{ਦਸ} \pa{ਵਜੇ} \pa{ਤੱਕ}}
\end{quote}

\begin{quote}
\emph{aṭṭh vaje tõ das vaje takk}
\end{quote}

\begin{quote}
from eight o'clock until ten o'clock
\end{quote}

The start is followed by \textbf{\pa{ਤੋਂ}}; the end limit is followed by \textbf{\pa{ਤੱਕ}}. Both are postpositions and each comes after the words it relates to. This step owns only the bounded from-until time frame.
\end{grammarlens}

\subsection*{Your turn}
\begin{itemize}
\raggedright
  \item \emph{Say it:} \textbf{\pa{ਅੱਠ} \pa{ਵਜੇ} \pa{ਤੋਂ}}
  \item \emph{Add:} \textbf{\pa{ਦਸ} \pa{ਵਜੇ} \pa{ਤੱਕ}}
  \item \emph{Read it:} \textbf{\pa{ਅੱਠ} \pa{ਵਜੇ} \pa{ਤੋਂ} \pa{ਦਸ} \pa{ਵਜੇ} \pa{ਤੱਕ}}
  \item \emph{Check:} start + \textbf{\pa{ਤੋਂ}}; end + \textbf{\pa{ਤੱਕ}}
\end{itemize}

\begin{tcolorbox}[breakable,colback=teal!4,colframe=teal!35!black,title={Before you move on}]
Give the whole span from eight until ten. (\textbf{\pa{ਅੱਠ} \pa{ਵਜੇ} \pa{ਤੋਂ} \pa{ਦਸ} \pa{ਵਜੇ} \pa{ਤੱਕ}}.) Which word follows the limit? (\textbf{\pa{ਤੱਕ}}.)

Source: Ranbir Johal, \emph{Speaking and Writing Punjabi} (2018), pp. 67–70, \href{https://openlibrary-repo.ecampusontario.ca/jspui/bitstream/123456789/492/2/2018Pun1100CoursePackage.pdf}{\pa{ਤੋਂ} … \pa{ਤੱਕ} and whole-hour schedule examples}.
\end{tcolorbox}
\section[\texorpdfstring{\pa{ਅਮਨ} \pa{ਨਾਲ} \pa{ਕੰਮ} \pa{ਕਰੋ} (Aman naal kamm karo)}{ਅਮਨ ਨਾਲ ਕੰਮ ਕਰੋ (Aman naal kamm karo)}]{\pa{ਅਮਨ} \pa{ਨਾਲ} \pa{ਕੰਮ} \pa{ਕਰੋ} — name one companion}

\label{lesson:PA-C178-naal}



\begin{quote}
Recall the practice name \textbf{\pa{ਅਮਨ}}, the familiar work word \textbf{\pa{ਕੰਮ}}, and polite \textbf{\pa{ਕਰੋ}}. Keep all three unchanged.
\end{quote}

\begin{grammarlens}[title={Grammar Lens: \pa{ਨਾਲ} follows the companion noun}]
\begin{quote}
\textbf{\pa{ਅਮਨ} \pa{ਨਾਲ} \pa{ਕੰਮ} \pa{ਕਰੋ।}}
\end{quote}

\begin{quote}
\emph{Aman nāl kamm karo.}
\end{quote}

\begin{quote}
Work with Aman.
\end{quote}

As with the last two forms, the postposition comes after its noun. The companion is \textbf{\pa{ਅਮਨ}}, so the pair is \textbf{\pa{ਅਮਨ} \pa{ਨਾਲ}}. This lesson teaches only the companion meaning of \textbf{\pa{ਨਾਲ}} in this polite instruction.
\end{grammarlens}

\subsection*{Your turn}
\begin{itemize}
\raggedright
  \item \emph{Hear:} \textbf{\pa{ਅਮਨ} \pa{ਨਾਲ}}; YOU SAY: with Aman
  \item \emph{Say it:} \textbf{\pa{ਅਮਨ} \pa{ਨਾਲ} \pa{ਕੰਮ} \pa{ਕਰੋ}}
  \item \emph{Read it:} \textbf{\pa{ਅਮਨ} \pa{ਨਾਲ} \pa{ਕੰਮ} \pa{ਕਰੋ}}
\end{itemize}

\begin{tcolorbox}[breakable,colback=teal!4,colframe=teal!35!black,title={Before you move on}]
Ask politely for work with Aman. (\textbf{\pa{ਅਮਨ} \pa{ਨਾਲ} \pa{ਕੰਮ} \pa{ਕਰੋ}}.) Which word follows the companion? (\textbf{\pa{ਨਾਲ}}.)

Source: Ranbir Johal, \emph{Speaking and Writing Punjabi} (2018), pp. 76 and 84, \href{https://openlibrary-repo.ecampusontario.ca/jspui/bitstream/123456789/492/2/2018Pun1100CoursePackage.pdf}{postpositions and \pa{ਨਾਲ}}.
\end{tcolorbox}
\section[\texorpdfstring{\pa{ਅਮਨ} \pa{ਨਾਲ} \pa{ਅੱਠ} \pa{ਵਜੇ} \pa{ਤੋਂ} … (Aman naal atth …)}{ਅਮਨ ਨਾਲ ਅੱਠ ਵਜੇ ਤੋਂ … (Aman naal atth …)}]{\pa{ਦਸ} \pa{ਵਜੇ} \pa{ਤੋਂ} \pa{ਬਾਅਦ} — place one step after a time}

\label{lesson:PA-C178-baad}



\begin{quote}
Say the schedule span \textbf{\pa{ਅੱਠ} \pa{ਵਜੇ} \pa{ਤੋਂ} \pa{ਦਸ} \pa{ਵਜੇ} \pa{ਤੱਕ}}. Then recall the companion pair \textbf{\pa{ਅਮਨ} \pa{ਨਾਲ}}. Nothing in those pieces changes.
\end{quote}

\begin{grammarlens}[title={Grammar Lens: \pa{ਤੋਂ} \pa{ਬਾਅਦ} marks the later step}]
\begin{quote}
\textbf{\pa{ਦਸ} \pa{ਵਜੇ} \pa{ਤੋਂ} \pa{ਬਾਅਦ}} — \emph{das vaje tõ bāad} — after ten o'clock
\end{quote}

Keep \textbf{\pa{ਤੋਂ}} before \textbf{\pa{ਬਾਅਦ}}. The time comes first, then the two-word sequence. This lesson owns only this bounded time use; it does not claim every temporal or clausal use of \textbf{\pa{ਬਾਅਦ}}.
\end{grammarlens}

\subsection*{Your turn: one two-sentence schedule}
\begin{quote}
\textbf{\pa{ਅਮਨ} \pa{ਨਾਲ} \pa{ਅੱਠ} \pa{ਵਜੇ} \pa{ਤੋਂ} \pa{ਦਸ} \pa{ਵਜੇ} \pa{ਤੱਕ} \pa{ਕੰਮ} \pa{ਕਰੋ।} \pa{ਦਸ} \pa{ਵਜੇ} \pa{ਤੋਂ} \pa{ਬਾਅਦ} \pa{ਘਰ} \pa{ਆਉਣਾ।}}
\end{quote}

\begin{itemize}
\raggedright
  \item \emph{Hear:} the two-sentence schedule; YOU NAME: companion, start, limit, later step
  \item \emph{Say it:} \textbf{\pa{ਅਮਨ} \pa{ਨਾਲ} \pa{ਅੱਠ} \pa{ਵਜੇ} \pa{ਤੋਂ} \pa{ਦਸ} \pa{ਵਜੇ} \pa{ਤੱਕ} \pa{ਕੰਮ} \pa{ਕਰੋ।} \pa{ਦਸ} \pa{ਵਜੇ} \pa{ਤੋਂ} \pa{ਬਾਅਦ} \pa{ਘਰ} \pa{ਆਉਣਾ।}}
  \item \emph{Read it:} \textbf{\pa{ਅਮਨ} \pa{ਨਾਲ} \pa{ਅੱਠ} \pa{ਵਜੇ} \pa{ਤੋਂ} \pa{ਦਸ} \pa{ਵਜੇ} \pa{ਤੱਕ} \pa{ਕੰਮ} \pa{ਕਰੋ।} \pa{ਦਸ} \pa{ਵਜੇ} \pa{ਤੋਂ} \pa{ਬਾਅਦ} \pa{ਘਰ} \pa{ਆਉਣਾ।}}
  \item \emph{Write it:} \textbf{\pa{ਅਮਨ} \pa{ਨਾਲ} \pa{ਅੱਠ} \pa{ਵਜੇ} \pa{ਤੋਂ} \pa{ਦਸ} \pa{ਵਜੇ} \pa{ਤੱਕ} \pa{ਕੰਮ} \pa{ਕਰੋ।} \pa{ਦਸ} \pa{ਵਜੇ} \pa{ਤੋਂ} \pa{ਬਾਅਦ} \pa{ਘਰ} \pa{ਆਉਣਾ।}} after covering the model
  \item \emph{Check:} \textbf{\pa{ਨਾਲ}} after the companion; \textbf{\pa{ਤੋਂ}} after the start; \textbf{\pa{ਤੱਕ}} after the limit; \textbf{\pa{ਤੋਂ} \pa{ਬਾਅਦ}} after the later time
\end{itemize}

\begin{tcolorbox}[breakable,colback=teal!4,colframe=teal!35!black,title={Before you move on}]
Rebuild the two-sentence schedule without looking. Then point to the five newly taught forms.

This is the first prerequisite tranche exposed by Punjabi A2 mock drafting. It does not unblock Punjabi A2 mock 1 by itself. It is not an A2 mock, not A2 readiness or attainment, and not external or human validation.

Sources: Ranbir Johal, \emph{Speaking and Writing Punjabi} (2018), pp. 67–70, 76, and 84–85, \href{https://openlibrary-repo.ecampusontario.ca/jspui/bitstream/123456789/492/2/2018Pun1100CoursePackage.pdf}{time and postposition frames}; \emph{Online Punjabi Teaching}, \href{https://www.learnpunjabi.org/Adverb.html}{clausal adverbs with \pa{ਤੋਂ} \pa{ਬਾਅਦ}}.
\end{tcolorbox}
